% Подробное методическое пояснение к лабораторной работе №1
% Компиляция: pdflatex guide.tex (дважды, из папки guide)
\documentclass[11pt,a4paper]{article}

\usepackage[T2A]{fontenc}
\usepackage[utf8]{inputenc}
\usepackage[russian]{babel}
\usepackage[left=2.2cm,right=1.8cm,top=2cm,bottom=2cm]{geometry}
\usepackage{booktabs}
\usepackage{longtable}
\usepackage{array}
\usepackage{graphicx}
\usepackage{enumitem}
\usepackage{amsmath}
\usepackage{xcolor}
\usepackage{mdframed}
\usepackage{caption}
\usepackage[hidelinks]{hyperref}
\usepackage[expansion=false]{microtype}

\captionsetup{font=small,labelfont=bf}
\setlength{\parindent}{1.25em}
\setlength{\emergencystretch}{3em}
\renewcommand{\arraystretch}{1.2}

\newcolumntype{L}[1]{>{\raggedright\arraybackslash}p{#1}}
\newcommand{\code}[1]{{\ttfamily #1}}

\definecolor{boxbg}{RGB}{242,242,247}
\definecolor{boxline}{RGB}{120,120,140}
\newmdenv[backgroundcolor=boxbg,linecolor=boxline,linewidth=0.6pt,
          innertopmargin=6pt,innerbottommargin=6pt,skipabove=8pt,skipbelow=8pt]{note}

\title{\textbf{Лабораторная работа №1: подробный разбор}\\[4pt]
\large Что мы делали, зачем и как это повторить\\[2pt]
\normalsize BCR::ABL1 тирозинкиназа --- WT и T315I}
\author{Шпока Владислав, Пясецкий Кирилл --- 4 курс, 3 группа}
\date{25 сентября 2026 г.}

\begin{document}
\sloppy
\maketitle

\begin{note}
\textbf{Зачем этот документ.} Отчёт (\code{report.pdf}) отвечает на вопросы задания
и содержит только результат. Здесь --- всё остальное: что означает каждый термин,
почему выбран тот или иной вариант, откуда взялось каждое число и что отвечать,
если на защите спросят «а почему именно так». Документ рассчитан на человека,
который открывает эту тему впервые.
\end{note}

\tableofcontents
\newpage

\section{Картина целиком: что вообще происходит в этой работе}

\subsection{Идея за заданием}

Разработка лекарства начинается не с молекулы, а с вопроса «на что воздействовать».
Нужно найти белок, чьё изменение вызывает болезнь, --- такой белок называют
\emph{молекулярной мишенью} (drug target). Дальше нужно убедиться, что по этому
белку вообще есть данные, с которыми можно работать на компьютере: известна ли
его пространственная структура, есть ли вещества, которые на него действуют, есть
ли измеренные числа, описывающие силу этого действия.

Вся работа №1 --- это именно такая проверка. Мы ничего не моделируем и не считаем
энергии: мы собираем «досье» на мишень и в конце честно отвечаем на вопрос
«данных достаточно или нет?». Всё собранное потом используется в следующих работах:
структуры --- для анализа сайта связывания и докинга, SMILES --- для хемоинформатики,
числа активности --- для сравнения и, возможно, построения моделей.

\subsection{Почему мишень --- BCR::ABL1}

BCR::ABL1 --- учебный «идеальный» пример. Для него известно буквально всё: причина
возникновения (одна конкретная хромосомная перестройка), болезнь (хронический
миелоидный лейкоз), десятки кристаллических структур, шесть одобренных препаратов
и тысячи измерений активности. Плюс у него есть знаменитая мутация устойчивости
T315I, на которой видно, как \emph{одна} замена аминокислоты обрушивает работу
лекарства. Это удобная пара для сравнения: «до» и «после».

\subsection{Из чего состоит результат}

\begin{center}
\begin{tabular}{L{4.6cm}L{10cm}}
\toprule
\textbf{Что сдаём} & \textbf{Что это} \\
\midrule
\code{report.pdf} & Отчёт по 13 пунктам задания со всеми таблицами \\
\code{BCR\_ABL\_dataset.csv} & Набор данных: 78 строк, одна строка --- одно измерение \\
\code{BCR\_ABL\_dataset.xlsx} & То же в Excel \\
\code{figures/fig\_activity.png} & График сравнения активности WT и T315I \\
\code{scripts/} & Скрипты, которыми всё выгружено \\
\code{data/} & Сырые ответы баз данных в JSON \\
\bottomrule
\end{tabular}
\end{center}

\section{Биология: минимум, который нужно понимать}

\subsection{Что такое киназа и почему она важна}

\textbf{Киназа} --- фермент, который переносит фосфатную группу с молекулы ATP на
другой белок. Это «выключатель»: у белка появляется заряженная фосфатная группа,
он меняет форму и начинает (или перестаёт) работать. Так клетка передаёт сигналы:
«расти», «делиться», «двигаться».

\textbf{Тирозинкиназа} --- киназа, которая фосфорилирует именно остатки тирозина.
ABL1 --- \emph{нерецепторная} тирозинкиназа: она не сидит в мембране, а работает
внутри клетки.

Ключевой момент: чтобы фосфорилировать что-то, киназе нужен ATP, и он связывается
в определённой щели между двумя «долями» белка. Эта щель --- \textbf{ATP-карман}.
Почти все лекарства против киназ --- это молекулы, которые занимают этот карман
вместо ATP. Отсюда термин \emph{ATP-конкурентный ингибитор}.

\begin{note}
\textbf{Интуиция.} Киназа --- это «раскладушка» из двух долей (N-доля и C-доля).
ATP садится в щель между ними. Лекарство --- «пробка», которая эту щель занимает.
Дальше вся история про T315I --- это история про то, что на входе в щель торчит
один остаток, и если он становится больше, пробка перестаёт входить.
\end{note}

\subsection{Нормальная ABL1 и её «предохранители»}

Здоровая ABL1 не работает постоянно --- иначе клетка бы бесконтрольно делилась.
У неё есть три механизма самоторможения (в UniProt это раздел
\emph{Activity regulation}):
\begin{enumerate}[leftmargin=*,topsep=2pt,itemsep=2pt]
\item на самом N-конце к белку пришит жирнокислотный «хвостик» (миристоил,
на остатке Gly2), который вставляется в карман на C-доле киназного домена и
фиксирует белок в неактивной форме;
\item домен SH3 прижимается к линкеру между SH2 и киназным доменом --- как
защёлка;
\item участок CAP (остатки 1--60) дополнительно стабилизирует эту «закрытую»
конструкцию.
\end{enumerate}
Когда клетке нужна активность ABL1, защёлки снимаются, петля активации
фосфорилируется по Tyr393, и фермент включается.

\subsection{Что ломает транслокация t(9;22)}

Ген \code{ABL1} лежит на хромосоме~9, ген \code{BCR} --- на хромосоме~22. Иногда
в клетке-предшественнице крови происходит обмен участками между этими хромосомами:
\textbf{реципрокная транслокация t(9;22)(q34;q11)}. Укороченная хромосома~22 ---
это и есть \emph{филадельфийская хромосома}. На ней оказывается химерный ген
\code{BCR::ABL1}.

Что получается на уровне белка:
\begin{itemize}[leftmargin=*,topsep=2pt,itemsep=2pt]
\item ABL1 теряет свой N-конец (остатки 1--26) --- а вместе с ним Gly2 и
миристоильный «предохранитель»;
\item вместо него приделывается кусок BCR, содержащий \emph{coiled-coil} домен ---
участок, который заставляет молекулы слипаться друг с другом (олигомеризоваться);
\item слипшиеся молекулы фосфорилируют друг друга по Tyr393 --- киназа включается
сама себя и больше не выключается.
\end{itemize}

Дальше --- цепочка: постоянный сигнал $\rightarrow$ постоянное фосфорилирование
белков-посредников (CRKL, STAT5 и др.) $\rightarrow$ активация путей роста
(RAS/MAPK, PI3K/AKT, JAK/STAT) $\rightarrow$ миелоидные клетки делятся без
внешней команды и не умирают $\rightarrow$ хронический миелоидный лейкоз.

\begin{note}
\textbf{Что важно для отчёта (пункт 2 задания):} преподаватель хочет увидеть не
пересказ учебника про лейкоз, а именно эту логическую цепочку из шести звеньев.
В отчёте она оформлена таблицей «Этап --- Содержание».
\end{note}

\subsection{Почему это хорошая мишень}

Три причины, которые стоит назвать:
\begin{enumerate}[leftmargin=*,topsep=2pt,itemsep=2pt]
\item \textbf{Драйверная мутация.} Не «один из факторов», а единственная причина
болезни. Подавили белок --- остановили болезнь.
\item \textbf{Избирательность.} Патологический белок есть только в опухолевых
клетках; нормальные клетки его не содержат, значит, у лекарства есть шанс быть
не слишком токсичным.
\item \textbf{«Друггабельность».} У белка есть готовый, хорошо очерченный карман
для малой молекулы (ATP-карман). Это редкость --- многие важные белки (например,
транскрипционные факторы) карманов не имеют, и подобрать к ним молекулу почти
невозможно.
\end{enumerate}

\subsection{T315I: почему одна буква решает всё}

\textbf{Обозначение.} T315I читается так: в положении 315 был Thr (треонин, T),
стал Ile (изолейцин, I). Аминокислоты обозначаются либо тремя буквами (Thr), либо
одной (T).

\textbf{Где этот остаток.} В самом конце «шарнира» (hinge) --- петли, соединяющей
две доли киназы, прямо на входе в ATP-карман. В UniProt сайт связывания ATP
указан как 316--322, то есть 315-й стоит вплотную перед ним.

\textbf{Почему «gatekeeper» (привратник).} За ATP-карманом есть дополнительная
полость --- \emph{back pocket}, или «задний карман». Многие ингибиторы (иматиниб,
нилотиниб) не помещаются в один только ATP-карман: часть их молекулы уходит в
заднюю полость. Пройдёт ли молекула туда, зависит от размера боковой цепи
остатка 315: он стоит ровно в проходе. Маленький треонин --- проход открыт;
крупный изолейцин --- закрыт.

\textbf{Два эффекта замены:}
\begin{enumerate}[leftmargin=*,topsep=2pt,itemsep=2pt]
\item \emph{Потеря водородной связи.} У треонина в боковой цепи есть группа --OH,
которая образует водородную связь с молекулой ингибитора. У изолейцина
(углеводородная боковая цепь) образовать такую связь нечем.
\item \emph{Стерическое препятствие.} Изолейцин объёмнее: примерно 167 \AA$^3$
против 116 \AA$^3$. Лишние атомы занимают место, где раньше лежала часть молекулы
ингибитора.
\end{enumerate}

Эффект второй причины сильнее. Потеря одной H-связи (4--8 кДж/моль) изменила бы
аффинность в несколько раз: связь между энергией и константой связывания ---
логарифмическая, $\Delta G = -RT\ln K$, и при комнатной температуре
$RT \approx 2{,}5$ кДж/моль, так что 6 кДж/моль дают примерно десятикратное
изменение. А наблюдаемое падение для дазатиниба --- в десятки тысяч раз. Значит,
дело не в энергетике одной связи, а в том, что молекула \emph{геометрически} не
помещается.

\textbf{Кто выживает при T315I:}
\begin{itemize}[leftmargin=*,topsep=2pt,itemsep=2pt]
\item \textbf{Понатиниб} --- в его молекуле есть линейный фрагмент с тройной
связью (\code{C\#C} в SMILES), тонкий как игла. Он проходит мимо увеличенного
изолейцина, не упираясь в него.
\item \textbf{Асциминиб} --- вообще не садится в ATP-карман. Он занимает тот
самый миристоильный карман, куда в здоровом белке вставляется жирный «хвостик».
Этот механизм назван STAMP (\emph{Specifically Targeting the ABL Myristoyl Pocket}).
Мутация в ATP-кармане ему просто безразлична.
\item \textbf{Аксетиниб} --- препарат, изначально разработанный против другой
киназы (VEGFR) и одобренный для рака почки. Выяснилось, что против T315I он
работает \emph{лучше}, чем против нормального ABL1. Это классический пример
\emph{перепрофилирования} (drug repurposing).
\end{itemize}

\section{Базы данных: что где лежит}

\begin{center}
\small
\begin{tabular}{L{2.4cm}L{5.4cm}L{6.6cm}}
\toprule
\textbf{База} & \textbf{Что хранит} & \textbf{Что мы взяли} \\
\midrule
UniProt & Белки: последовательность, домены, функции, болезни & Запись P00519: длина 1130 а.о., киназный домен 242--493, сайты ATP, точка разрыва 26/27 \\
RCSB PDB & Экспериментальные 3D-структуры (координаты атомов) & 109 записей, связанных с ABL1; 68 с киназным доменом; выбрали 2HYY и 3QRJ \\
AlphaFold DB & Предсказанные нейросетью структуры & Модель AF-P00519-F1, покрывает остатки 1--1130 \\
ChEMBL & Биоактивности из статей: IC$_{50}$, K$_i$, K$_d$ & 78 измерений для 7 соединений против WT и T315I \\
BindingDB & То же, с уклоном в данные связывания & Использована как перекрёстная проверка \\
PubChem & Химические вещества: структуры, идентификаторы & CID, формула, масса, SMILES для 7 соединений \\
\bottomrule
\end{tabular}
\end{center}

\subsection{Что такое «reviewed» в UniProt}

UniProt состоит из двух частей. \textbf{Swiss-Prot} (помечено \emph{reviewed},
значок с галочкой) --- записи, которые вручную проверил куратор: функции
подтверждены ссылками на статьи, домены размечены. \textbf{TrEMBL}
(\emph{unreviewed}) --- автоматически сгенерированные записи, их на порядки больше
и качество ниже. Задание требует именно reviewed-запись, поэтому мы берём
\code{ABL1\_HUMAN} / \code{P00519}.

\subsection{Что такое запись PDB и что в ней смотреть}

Одна запись PDB --- это результат \emph{одного эксперимента}: конкретный белок
(часто только его фрагмент), возможно с конкретным лигандом, закристаллизованный
в конкретных условиях. Поэтому у одного белка бывают десятки записей. Смотреть
нужно на:

\begin{center}
\small
\begin{tabular}{L{4.0cm}L{10.6cm}}
\toprule
\textbf{Поле} & \textbf{Что означает} \\
\midrule
Experimental Method & X-ray diffraction (рентген), NMR (ЯМР), cryo-EM. Для нашей задачи лучше рентген: даёт точные координаты и лиганд \\
Resolution & Разрешение в ангстремах. \emph{Меньше --- лучше.} До 2.0 \AA\ --- хорошо, 2.0--2.5 --- приемлемо, хуже 3.0 --- положение отдельных боковых цепей уже ненадёжно \\
Ligands & Малые молекулы в структуре. Отличайте лекарство от «мусора»: SO4, GOL (глицерин), CL, NA --- это соли и криопротектор из буфера, а не лиганд \\
Macromolecules / Entity & Какой кусок белка есть в структуре (например, 229--500 --- только киназный домен) \\
Mutation(s) & Введённые замены. Бывает, что кроме нужной T315I есть ещё лишние --- такие структуры хуже \\
\bottomrule
\end{tabular}
\end{center}

\begin{note}
\textbf{Ловушка, на которую легко попасть.} В поле «Mutation» часть записей
пишет \code{YES} без расшифровки (так в 2V7A), а часть использует нумерацию другой
изоформы: в \code{5MO4} стоит \code{T334I}, и это \emph{та же самая} мутация
T315I, просто отсчёт начинается на 19 остатков раньше. Поэтому в работе мы
определяли форму белка не по этому полю, а по самой последовательности ---
см.~раздел~\ref{sec:howto}.
\end{note}

\subsection{IC$_{50}$, K$_i$, K$_d$: в чём разница}

Это три разных числа, и путать их нельзя.

\begin{center}
\small
\begin{tabular}{L{1.5cm}L{5.2cm}L{7.8cm}}
\toprule
\textbf{Вел.} & \textbf{Определение} & \textbf{Особенности} \\
\midrule
K$_d$ & Константа диссоциации комплекса белок--лиганд & Чистая термодинамика: насколько прочно молекула держится. Не зависит от ATP \\
K$_i$ & Константа ингибирования & Похожа на K$_d$, но получена из кинетики фермента \\
IC$_{50}$ & Концентрация, при которой активность фермента падает вдвое & Зависит от условий: концентрации ATP, времени инкубации, субстрата. В разных лабораториях будет разной \\
\bottomrule
\end{tabular}
\end{center}

Во всех трёх случаях \emph{меньше --- значит сильнее}. 1~нМ --- очень сильный
ингибитор, 1000~нМ (=\,1~мкМ) --- слабый.

Связь IC$_{50}$ и K$_i$ для конкурентного ингибитора даёт уравнение
Ченга--Прусоффа:
\[
K_i = \frac{\mathrm{IC}_{50}}{1 + [\mathrm{ATP}]/K_m}.
\]
Отсюда видно, почему две лаборатории получают разные IC$_{50}$: если у одной в
пробе 10~мкМ ATP, а у другой 1~мМ, числа разойдутся в десятки раз при одинаковой
реальной силе связывания.

\begin{note}
\textbf{Это --- ответ на вопрос в пункте~8 задания} («можно ли сравнивать все
найденные значения напрямую?»). Нет, нельзя: разные величины, разные условия,
разные конструкции белка. Сравнивать можно только внутри одной публикации и
одного типа анализа. В нашем наборе IC$_{50}$ иматиниба против WT меняется от 75
до 7700 нМ --- сто раз разницы, и это при одном и том же веществе и одной и той
же мишени.
\end{note}

\subsection{Что такое «цензурированное значение»}

В таблицах встречается \code{>10000 нМ}. Это не измеренная величина, а граница:
«в исследованном диапазоне концентраций (до 10 мкМ) подавления не увидели».
Настоящее значение может быть 20 мкМ, а может 500 мкМ --- мы не знаем. Такие
данные:
\begin{itemize}[leftmargin=*,topsep=2pt,itemsep=1pt]
\item \emph{нельзя} подставлять в регрессию как число;
\item \emph{нужно} сохранять со знаком «$>$», а не молча превращать в 10000.
\end{itemize}
Именно поэтому в нашем датасете в столбце \code{Activity\_value} стоит строка
\code{>10000}, а не число.

\subsection{Что такое SMILES}

SMILES (\emph{Simplified Molecular Input Line Entry System}) --- способ записать
молекулу одной строкой. Читается примерно так:

\begin{center}
\small
\begin{tabular}{L{3.2cm}L{11.4cm}}
\toprule
\textbf{Фрагмент} & \textbf{Значение} \\
\midrule
\code{C}, \code{N}, \code{O} & атомы углерода, азота, кислорода (водороды подразумеваются) \\
\code{c1ccccc1} & ароматическое кольцо (строчные буквы = ароматика) \\
\code{C1CCCCC1} & обычный цикл: цифры отмечают, где кольцо замыкается \\
\code{=} \code{\#} & двойная и тройная связь \\
\code{( )} & ответвление \\
\code{/C=C/} & конфигурация двойной связи ($E$ или $Z$) \\
\code{[C@@H]} & стереоцентр (хиральный атом) с указанием ориентации \\
\bottomrule
\end{tabular}
\end{center}

Например, у понатиниба в SMILES есть \code{C\#CC4=CN=C5N4N=CC=C5} --- это та самая
«игла» с тройной связью, которая позволяет ему обходить мутировавший gatekeeper.
То есть по строке SMILES уже можно кое-что сказать о механизме.

\textbf{Почему SMILES нужен.} Компьютер не умеет работать с названием «иматиниб» ---
это просто слово. А SMILES --- это полное описание молекулярного графа: любая
библиотека (RDKit, Open Babel) прочитает его, построит 2D- или 3D-структуру,
посчитает дескрипторы и отпечатки, оценит сходство молекул. И всё это помещается
в одну ячейку таблицы --- в отличие от файла координат.

\section{Как именно всё было сделано}
\label{sec:howto}

Этот раздел --- воспроизводимая инструкция. Все шаги можно повторить как руками
через сайты, так и скриптами из папки \code{scripts/}.

\subsection{Шаг 1. UniProt}

\textbf{Руками:} зайти на \url{https://www.uniprot.org}, ввести \emph{ABL1 human},
выбрать запись с пометкой \emph{Reviewed} и названием \code{ABL1\_HUMAN}. Нужные
поля находятся в разделах \emph{Names \& Taxonomy}, \emph{Function},
\emph{Family \& Domains}, \emph{Sequence}.

\textbf{Скриптом:} \code{python scripts/fetch\_uniprot.py}. Скрипт обращается к
\code{https://rest.uniprot.org/uniprotkb/P00519.json} и печатает всё разом.

\textbf{Что получили:}
\begin{itemize}[leftmargin=*,topsep=2pt,itemsep=1pt]
\item длина 1130 а.о.; киназный домен 242--493; SH3 61--121; SH2 127--217;
\item сайты связывания ATP: 248--256, 271, 316--322; каталитический Asp363;
\item точка разрыва при транслокации: между 26 и 27 остатком (прямо в записи!);
\item остаток 315 --- действительно Thr, окружение \code{TREPPFYIITEFMTYGNLLD}.
\end{itemize}

\begin{note}
\textbf{Подводный камень: нумерация.} У ABL1 две изоформы. Изоформа 1a --- 1130
остатков, gatekeeper = Thr\textbf{315}. Изоформа 1b на 19 остатков длиннее, там
тот же остаток --- Thr\textbf{334}. В литературе и в PDB используются обе. Мы
везде работаем в нумерации 1a (как в UniProt P00519-1) и оговариваем это в отчёте.
\end{note}

\subsection{Шаг 2. Поиск структур в PDB}

\textbf{Руками:} на \url{https://www.rcsb.org} искать \emph{ABL1 kinase} и
\emph{ABL1 T315I}, затем открывать записи и смотреть поля из таблицы выше.

\textbf{Скриптом:} \code{python scripts/survey\_pdb.py}. Логика:
\begin{enumerate}[leftmargin=*,topsep=2pt,itemsep=2pt]
\item через Search API найти все записи, ссылающиеся на UniProt \code{P00519}
(человек) и \code{P00520} (мышь) --- получилось 109;
\item через Data API (GraphQL) забрать по ним метаданные, включая
последовательность;
\item определить форму белка \emph{по последовательности}: если в ней есть мотив
\code{FYIITEF} --- это WT, если \code{FYIIIEF} --- T315I. Третья буква I вместо T ---
это и есть мутация.
\end{enumerate}

\textbf{Почему именно так, а не по полю «Mutation».} Как уже сказано, поле
заполнено непоследовательно: где-то \code{YES}, где-то \code{T334I}. А
последовательность есть всегда и не врёт. Заодно этот же критерий отсеивает
записи, где киназного домена вообще нет (структуры отдельных SH3/SH2-доменов или
комплексов с ДНК): если мотива нет ни в каком виде, значит остаток 315 в
структуру не попал.

\textbf{Результат:} 68 записей содержат киназный домен; из них 60 --- WT, 8 ---
T315I; 64 из 68 содержат связанную малую молекулу; 65 получены рентгеном, 3 --- ЯМР.
Эти числа попали в отчёт и в итоговый вывод.

\subsection{Шаг 3. Выбор двух структур}

Критерии из задания и как мы их применили:

\begin{center}
\small
\begin{tabular}{L{5.2cm}L{9.4cm}}
\toprule
\textbf{Критерий} & \textbf{Как проверяли} \\
\midrule
Содержит ABL kinase domain & Проверено по выравниванию: 2HYY --- остатки 228--500, 3QRJ --- 229--499 \\
Есть известный ингибитор & 2HYY: STI (иматиниб); 3QRJ: 919 (ребастиниб) \\
Хорошие экспериментальные характеристики & 2HYY --- 2.40 \AA, 3QRJ --- 1.82 \AA\ (лучшее разрешение среди всех структур T315I) \\
Нужная форма белка & 2HYY --- Thr315; 3QRJ --- Ile315 (по последовательности) \\
Нет лишних мутаций & Обе записи «чистые». Для сравнения: 5MO4 содержит ещё и D363N (инактивация катализа), 3DK7 --- L445P \\
Полный сайт связывания & В обеих структурах лиганд разрешён целиком \\
Человеческий белок & Обе --- \emph{Homo sapiens}. Это важно: 1IEP, 3IK3, 3OY3, 2Z60 --- мышиные \\
\bottomrule
\end{tabular}
\end{center}

\textbf{Итог: 2HYY (WT) и 3QRJ (T315I).}

\textbf{Что можно ответить, если спросят про альтернативу.} Пара 4WA9/4TWP тоже
хороша --- это WT и T315I с одним и тем же лигандом (аксетинибом), полученные в
одной работе. Её плюс: любые различия между структурами почти наверняка вызваны
мутацией, а не разными лигандами. Минус: разрешение хуже (2.20 и 2.40 \AA), и
2HYY --- гораздо более цитируемый эталон для ATP-кармана. Мы выбрали 2HYY/3QRJ,
а 4WA9/4TWP держим как контрольную пару.

\subsection{Шаг 4. AlphaFold}

Открывается напрямую: \url{https://alphafold.ebi.ac.uk/entry/P00519}. Модель
\code{AF-P00519-F1}, покрывает всю последовательность 1--1130.

Что тут важно понять --- разницу трёх типов данных:
\begin{center}
\small
\begin{tabular}{L{3.0cm}L{5.4cm}L{6.2cm}}
\toprule
& \textbf{Что это} & \textbf{Чего там нет} \\
\midrule
Sequence (UniProt) & Порядок аминокислот, аннотации & Нет координат вообще \\
Structure (PDB) & Реально измеренные координаты атомов & Обычно только фрагмент белка \\
Prediction (AlphaFold) & Предсказанные координаты для всей цепи & Нет лигандов, ионов, воды; нет разрешения; одна усреднённая конформация \\
\bottomrule
\end{tabular}
\end{center}

Готовой модели \emph{BCR::ABL1 T315I} в AlphaFold DB нет и быть не может: база
строится по референсным протеомам, а слитный белок --- продукт соматической
перестройки в опухолевой клетке, его нет в «нормальном» геноме. Мутантную модель
при необходимости строят сами.

\subsection{Шаг 5. Ингибиторы и их активности}

\textbf{Руками:} в ChEMBL найти соединение (например, \emph{imatinib}), открыть
вкладку \emph{Activities}, отфильтровать по мишени \emph{ABL1} и по типу
показателя.

\textbf{Скриптом:} \code{python scripts/fetch\_chembl.py}. Запрос идёт к
\code{https://www.ebi.ac.uk/chembl/api/data/activity} с параметрами
\code{molecule\_chembl\_id}, \code{target\_chembl\_id} и
\code{standard\_type\_\_in=IC50,Ki,Kd}.

Две мишени в ChEMBL:
\begin{itemize}[leftmargin=*,topsep=2pt,itemsep=1pt]
\item \code{CHEMBL1862} --- ABL1, изолированная киназа;
\item \code{CHEMBL2096618} --- BCR-ABL1, слитный белок.
\end{itemize}
Форму (WT или мутант) даёт поле \code{assay\_variant\_mutation}: пусто --- WT,
\code{T315I} --- мутант. Есть и другие мутанты (E255K, Y253H, F317L\ldots) ---
их мы отбросили, они не входят в задание.

\textbf{Почему выбрали именно эти семь соединений.} Нужно было показать разное
поведение относительно T315I:
\begin{itemize}[leftmargin=*,topsep=2pt,itemsep=1pt]
\item иматиниб, нилотиниб, дазатиниб --- теряют активность полностью;
\item бозутиниб --- теряет частично (интересный промежуточный случай);
\item понатиниб --- сохраняет (ATP-конкурентный, но «обходит» gatekeeper);
\item асциминиб --- сохраняет (другой сайт связывания);
\item аксетиниб --- парадокс: против мутанта активнее, чем против нормы.
\end{itemize}

\textbf{Честный пробел.} Для асциминиба в ChEMBL не нашлось записи с явной
пометкой \code{T315I}, хотя из литературы известно, что он против этого мутанта
работает. В отчёте так и написано: «не найдено», с указанием, где искали. Задание
это прямо разрешает, и это правильнее, чем подставить число «по памяти».

\subsection{Шаг 6. PubChem}

\code{python scripts/fetch\_pubchem.py} --- запрос к PUG REST API по названию
вещества, свойства \code{MolecularFormula}, \code{MolecularWeight},
\code{SMILES}, \code{InChIKey}.

Для асциминиба и аксетиниба важно брать \emph{изомерный} SMILES: у первого есть
стереоцентр (\code{[C@@H]}), у второго --- $E$-конфигурация двойной связи
(\code{/C=C/}). Если взять запись без стереохимии, это будет уже, строго говоря,
другое вещество.

\subsection{Шаг 7. Сборка датасета}

\code{python scripts/build\_dataset.py} объединяет выгрузки в таблицу с колонками
из задания. Ключевые решения:

\begin{center}
\small
\begin{tabular}{L{4.2cm}L{10.4cm}}
\toprule
\textbf{Решение} & \textbf{Почему} \\
\midrule
Одна строка --- одно измерение & Так требует задание, и так правильно: усреднять значения из разных анализов нельзя (см. Ченга--Прусоффа) \\
Только отобранные публикации & Иначе в набор попадают случайные клеточные анализы без привязки к форме мишени \\
Знак «$>$» сохраняется & Цензурированные данные остаются видимыми \\
\code{PDB\_ID} где возможно & Связывает число с конкретной структурой; если такой структуры нет --- поле пустое, а не выдуманное \\
\code{Source} с PMID & Любое число можно проверить по первоисточнику \\
Разделитель «;», UTF-8 с BOM & Чтобы русский Excel открыл файл сразу и по столбцам \\
\bottomrule
\end{tabular}
\end{center}

Получилось 78 строк: 42 для WT и 36 для T315I, 7 соединений, 11 первоисточников.

\subsection{Шаг 8. График}

\code{python scripts/plot\_activity.py} строит столбчатую диаграмму: по одному
столбцу WT и T315I на каждое соединение.

\textbf{Три решения, которые стоит уметь объяснить:}
\begin{enumerate}[leftmargin=*,topsep=2pt,itemsep=2pt]
\item \textbf{Логарифмическая шкала по оси Y.} Значения различаются в $10^4$ раз.
На линейной шкале все столбцы, кроме дазатиниба против T315I, слились бы в одну
линию у нуля.
\item \textbf{Медиана, а не среднее.} Измерений на соединение несколько, и они
из разных статей; одно выбивающееся значение (7700 нМ для иматиниба) сдвинуло бы
среднее, а медиана устойчива.
\item \textbf{Знак «$>$» над столбцом.} Если среди измерений есть цензурированные,
столбец --- это заниженная оценка. Не пометить это --- значит соврать в графике.
\end{enumerate}

\textbf{Что видно на графике.} Иматиниб, нилотиниб, дазатиниб --- огромный скачок
вверх (потеря активности). Бозутиниб --- умеренный. Понатиниб --- почти ничего.
Аксетиниб --- столбец T315I \emph{ниже} столбца WT. Вывод в одну фразу: T315I --- не
«общее ухудшение связывания», а избирательный эффект, зависящий от того, опирается
ли молекула на контакт с gatekeeper-остатком.

\section{Что отвечать на защите}

\subsection{Вопросы, которые почти наверняка зададут}

\begin{description}[leftmargin=0pt,style=unboxed,itemsep=6pt]
\item[\textnormal{\textbf{Почему вы работаете с ABL1, а не с BCR::ABL1?}}]
Потому что патологическая активность целиком определяется киназным доменом ABL1,
все ингибиторы связываются в нём, мутация T315I --- в нём же, и структурных данных
по полноразмерному слитному белку просто не существует.

\item[\textnormal{\textbf{Что означает WT в этой работе?}}]
Форма киназного домена \emph{без мутации T315I}. Это не значит «нормальный
физиологический белок»: BCR::ABL1 патологичен в любом случае. WT здесь --- условное
обозначение из задания.

\item[\textnormal{\textbf{Почему выбрали 2HYY и 3QRJ?}}]
Обе --- человеческий киназный домен без посторонних мутаций, обе с ингибитором в
структуре, у 3QRJ лучшее разрешение среди всех T315I-структур (1.82 \AA), 2HYY ---
канонический эталон комплекса с иматинибом.

\item[\textnormal{\textbf{Почему у вас в таблице стоит «$>$10000», а не число?}}]
Это цензурированное значение: эксперимент показал только, что в исследованном
диапазоне подавления нет. Подставлять сюда 10000 как измеренную величину было бы
ошибкой.

\item[\textnormal{\textbf{Можно ли сравнивать IC$_{50}$ из разных статей?}}]
Нет. IC$_{50}$ зависит от концентрации ATP (Ченг--Прусофф), формата анализа и
конструкции белка. В нашем наборе одно и то же вещество против одной и той же
мишени даёт разброс в 100 раз. Сравнивать можно только внутри однородной группы.

\item[\textnormal{\textbf{Почему у аксетиниба активность против мутанта выше?}}]
Он связывается в конформации, которая не требует прохода за gatekeeper, и,
по-видимому, лучше подходит именно к геометрии мутантного кармана. Это
экспериментальный факт из работы 2015 года (PMID 25686603), подтверждённый
структурами 4WA9 и 4TWP.

\item[\textnormal{\textbf{Чем асциминиб отличается от остальных?}}]
Он не ATP-конкурентный. Он занимает миристоильный карман --- то место, куда в
здоровом белке вставляется N-концевой жирный «хвостик», и восстанавливает
автоингибирование. Поэтому мутации в ATP-кармане на него не влияют.

\item[\textnormal{\textbf{Достаточно ли данных для дальнейшей работы?}}]
Да, но с оговорками: структур T315I в 7{,}5 раза меньше, чем WT; часть данных
активности цензурирована; значения из разных источников несопоставимы напрямую;
модели слитного белка нет.
\end{description}

\subsection{Термины одной строкой}

\begin{center}
\small
\begin{longtable}{L{3.7cm}L{10.9cm}}
\toprule
\textbf{Термин} & \textbf{Значение} \\
\midrule
\endfirsthead
\toprule
\textbf{Термин} & \textbf{Значение} \\
\midrule
\endhead
Drug target & Молекула (обычно белок), на которую направлено действие лекарства \\
Киназа & Фермент, переносящий фосфат с ATP на белок-мишень \\
Kinase domain & Каталитическая часть киназы; у ABL1 --- остатки 242--493 \\
ATP-карман & Щель между долями киназы, где связывается ATP; сюда садятся ингибиторы \\
Gatekeeper & Остаток на входе в заднюю полость кармана; у ABL1 это Thr315 \\
Back pocket & Дополнительная гидрофобная полость за gatekeeper \\
DFG-in / DFG-out & Две конформации петли активации (по мотиву Asp-Phe-Gly): активная и неактивная. Ингибиторы типа I связывают DFG-in, типа II --- DFG-out \\
Петля активации & Участок 381--405; фосфорилирование Tyr393 включает киназу \\
Миристоильный карман & Полость на C-доле, куда вставляется N-концевой жирный «хвостик»; сайт связывания асциминиба \\
Fusion gene & Химерный ген из двух разных генов после хромосомной перестройки \\
t(9;22) & Транслокация, дающая филадельфийскую хромосому и ген BCR::ABL1 \\
TKI & Tyrosine Kinase Inhibitor --- ингибитор тирозинкиназы \\
IC$_{50}$ & Концентрация, снижающая активность фермента вдвое \\
K$_d$ / K$_i$ & Константы диссоциации комплекса / ингибирования \\
Цензурированное значение & Граница вида «$>$10000», а не измеренное число \\
SMILES & Строковая запись структуры молекулы \\
InChIKey & Хешированный идентификатор молекулы фиксированной длины \\
PDB ID & Четырёхсимвольный код записи в Protein Data Bank (например, 2HYY) \\
Resolution & Разрешение структуры в ангстремах; меньше --- лучше \\
Apo / holo & Структура без лиганда / с лигандом \\
pLDDT, PAE & Оценки достоверности предсказания AlphaFold (разбираются в другой работе) \\
Редокинг & Проверка протокола докинга: лиганд вынимают из структуры и пытаются вернуть на место \\
\bottomrule
\end{longtable}
\end{center}

\section{Как воспроизвести работу с нуля}

\subsection{Что нужно установить}

\begin{center}
\small
\begin{tabular}{L{4.6cm}L{10.0cm}}
\toprule
\textbf{Компонент} & \textbf{Зачем} \\
\midrule
Python 3.8+ & Запуск скриптов; хватает стандартной библиотеки \\
\code{matplotlib}, \code{numpy} & Построение графика: \code{pip install matplotlib numpy} \\
\code{openpyxl} (необязательно) & Сохранение .xlsx; без него делается только .csv \\
MiKTeX / TeX Live & Сборка отчёта в PDF \\
\bottomrule
\end{tabular}
\end{center}

\subsection{Порядок запуска}

\begin{center}
\small
\begin{tabular}{L{1.0cm}L{5.4cm}L{7.8cm}}
\toprule
\textbf{№} & \textbf{Команда} & \textbf{Что получится} \\
\midrule
1 & \code{python scripts/fetch\_uniprot.py} & \code{data/uniprot\_P00519.json}, данные для таблицы раздела 1 \\
2 & \code{python scripts/survey\_pdb.py} & \code{data/abl\_rows.json}, статистика по структурам \\
3 & \code{python scripts/fetch\_chembl.py} & \code{data/chembl\_activities.json} \\
4 & \code{python scripts/fetch\_pubchem.py} & \code{data/pubchem.json} \\
5 & \code{python scripts/build\_dataset.py} & \code{BCR\_ABL\_dataset.csv} и \code{.xlsx} \\
6 & \code{python scripts/plot\_activity.py} & \code{figures/fig\_activity.png} \\
7 & \code{cd report \&\& pdflatex report.tex} (дважды) & \code{report.pdf} \\
\bottomrule
\end{tabular}
\end{center}

\begin{note}
\textbf{Почему pdflatex запускают дважды.} При первом проходе LaTeX собирает
номера разделов, ссылок и страниц во вспомогательный файл \code{.aux}, при втором ---
подставляет их в текст. После одного прохода вместо номеров разделов будут
вопросительные знаки.
\end{note}

\subsection{Если что-то не работает}

\begin{center}
\small
\begin{tabular}{L{5.4cm}L{9.2cm}}
\toprule
\textbf{Симптом} & \textbf{Что делать} \\
\midrule
\code{socket.timeout} при обращении к ChEMBL & Сервер отвечает медленно; в скрипте уже есть повторы. Запустить ещё раз \\
В Excel вся строка в одной ячейке & Файл открыт с неверным разделителем. Данные $\rightarrow$ Текст по столбцам $\rightarrow$ разделитель «;». Или открыть .xlsx \\
Кракозябры вместо русского в CSV & Нужна кодировка UTF-8; файл уже сохранён с BOM, но старые версии Excel всё равно просят указать вручную \\
\code{! LaTeX Error: File ... not found} & MiKTeX докачивает пакеты при первой сборке --- разрешить установку и пересобрать \\
Вместо русского в PDF пустоты & Не подключены \code{fontenc[T2A]} и \code{babel[russian]} --- в наших файлах они есть, проверить преамбулу \\
\bottomrule
\end{tabular}
\end{center}

\section{Что дальше по курсу}

Собранное здесь --- вход для следующих работ:

\begin{center}
\small
\begin{tabular}{L{5.2cm}L{9.4cm}}
\toprule
\textbf{Что понадобится} & \textbf{Где будет использоваться} \\
\midrule
PDB 2HYY и 3QRJ & Анализ структуры белка, подготовка рецептора, определение binding site, докинг \\
SMILES семи соединений & Хемоинформатика: дескрипторы, отпечатки, сходство; подготовка лигандов \\
\code{BCR\_ABL\_dataset.csv} & Сравнение расчётных и экспериментальных величин; возможно, QSAR \\
Понимание T315I & Объяснение, почему докинг иматиниба в 3QRJ должен давать худший результат, чем в 2HYY \\
AlphaFold-модель & Отдельная работа про pLDDT, PAE и надёжность предсказаний \\
\bottomrule
\end{tabular}
\end{center}

\medskip
\noindent Главное, что стоит запомнить из этой работы: прежде чем что-то считать,
нужно понять, из чего сделаны исходные данные. Числа из баз --- не абсолютная
истина: у них есть метод получения, условия эксперимента, единицы измерения и
ограничения. Работа №1 --- как раз про то, чтобы научиться это видеть.

\end{document}
